\documentclass{beamer}

%================ ENCODING =================
\usepackage[utf8]{inputenc}
\usepackage[T1]{fontenc}
\usepackage[vietnamese]{babel}

%================ PACKAGES =================
\usepackage{tikz}
\usetikzlibrary{arrows.meta,positioning}
\usepackage[most]{tcolorbox}
\usepackage{amsmath}
\usepackage{xcolor}
\usepackage{fontawesome5}

%================ COLORS (PASTEL LIGHT) =================
\definecolor{softblue}{RGB}{173,216,255}
\definecolor{softpurple}{RGB}{210,180,255}
\definecolor{bg1}{RGB}{245,245,255}
\definecolor{bg2}{RGB}{230,220,255}

%================ BACKGROUND =================
\usebackgroundtemplate{
\begin{tikzpicture}[remember picture,overlay]
\shade[left color=bg1,right color=bg2]
(current page.south west) rectangle (current page.north east);

\fill[softblue,opacity=0.2] (current page.north east) circle (2cm);
\fill[softpurple,opacity=0.2] (current page.south west) circle (2cm);
\end{tikzpicture}
}

%================ STYLE =================
\setbeamercolor{normal text}{fg=black}
\setbeamertemplate{navigation symbols}{}

%================ FOOTLINE =================
\setbeamercolor{footline}{fg=black,bg=white}
\setbeamertemplate{footline}{
\leavevmode
\hbox{
\begin{beamercolorbox}[wd=0.5\paperwidth,ht=2.5ex,dp=1ex,left]{footline}
\hspace{0.3cm} Hoàng Tuệ Anh
\end{beamercolorbox}
\begin{beamercolorbox}[wd=0.5\paperwidth,ht=2.5ex,dp=1ex,right]{footline}
\insertframenumber{} / \inserttotalframenumber \hspace{0.3cm}
\end{beamercolorbox}
}
}

%================ BOX =================
\tcbset{
colback=white,
colframe=softblue!70!black,
arc=10pt,
boxrule=0.8pt
}

\begin{document}

%================ SLIDE 1 =================
\begin{frame}{\faProjectDiagram\ Biểu diễn đồ thị}
\begin{tcolorbox}
\begin{itemize}
\item Đồ thị $G = (V, E)$ với $V$ là tập đỉnh, $E$ là tập cạnh
\item Đồ thị có thể vô hướng hoặc có hướng
\item Biểu diễn trực quan bằng hình học (graph drawing)
\item Mỗi cạnh nối hai đỉnh biểu diễn quan hệ
\item Là mô hình cơ bản trong toán rời rạc
\end{itemize}
\end{tcolorbox}

\begin{center}
\begin{tikzpicture}[scale=1]
\node[circle,draw,fill=softblue!60] (1) at (0,0) {1};
\node[circle,draw,fill=softblue!60] (2) at (2,0) {2};
\node[circle,draw,fill=softpurple!60] (3) at (1,1.5) {3};
\node[circle,draw,fill=softpurple!60] (4) at (1,-1.5) {4};

\draw (1)--(2);
\draw (1)--(3);
\draw (2)--(4);
\end{tikzpicture}
\end{center}
\end{frame}

%================ SLIDE 2 =================
\begin{frame}{\faList\ Danh sách kề (Adjacency List)}
\begin{tcolorbox}
\begin{itemize}
\item Mỗi đỉnh lưu danh sách các đỉnh kề
\item Thường cài đặt bằng vector hoặc linked list
\item Tiết kiệm bộ nhớ với đồ thị thưa
\item Độ phức tạp lưu trữ: $O(V + E)$
\item Thuận lợi cho duyệt DFS/BFS
\end{itemize}
\end{tcolorbox}

\begin{center}
1: 2,3 \\
2: 1,4 \\
3: 1 \\
4: 2
\end{center}
\end{frame}

%================ SLIDE: MINH HỌA DANH SÁCH KỀ =================
\begin{frame}{\faProjectDiagram\ Minh họa danh sách kề}
\begin{tcolorbox}
\begin{itemize}
\item Đồ thị gồm 4 đỉnh: 1, 2, 3, 4
\item Các cạnh: (1,2), (1,3), (2,4)
\item Biểu diễn bằng danh sách kề
\item Mỗi đỉnh lưu các đỉnh kề trực tiếp
\item Dễ dùng cho DFS và BFS
\end{itemize}
\end{tcolorbox}

\vspace{0.3cm}

\begin{columns}
% Cột trái: đồ thị
\column{0.5\textwidth}
\centering
\begin{tikzpicture}[scale=1]

\node[circle,draw,fill=softblue!60] (1) at (0,0) {1};
\node[circle,draw,fill=softblue!60] (2) at (2,0) {2};
\node[circle,draw,fill=softpurple!60] (3) at (1,1.5) {3};
\node[circle,draw,fill=softpurple!60] (4) at (1,-1.5) {4};

\draw (1)--(2);
\draw (1)--(3);
\draw (2)--(4);

\end{tikzpicture}

% Cột phải: danh sách kề
\column{0.5\textwidth}
\begin{tcolorbox}[colframe=softpurple!70!black]
\textbf{Danh sách kề:}
\begin{itemize}
\item 1: 2, 3
\item 2: 1, 4
\item 3: 1
\item 4: 2
\end{itemize}
\end{tcolorbox}

\end{columns}

\end{frame}

%================ SLIDE 3 =================
\begin{frame}{\faTable\ Ma trận kề (Adjacency Matrix)}
\begin{tcolorbox}
\begin{itemize}
\item Ma trận $A[i][j]$ biểu diễn cạnh giữa $i, j$
\item $A[i][j] = 1$ nếu có cạnh, ngược lại = 0
\item Phù hợp với đồ thị dày
\item Truy cập cạnh trong $O(1)$
\item Tốn bộ nhớ $O(V^2)$
\end{itemize}
\end{tcolorbox}

\[
\begin{bmatrix}
0 & 1 & 1 & 0 \\
1 & 0 & 0 & 1 \\
1 & 0 & 0 & 0 \\
0 & 1 & 0 & 0
\end{bmatrix}
\]
\end{frame}

%================ SLIDE: MINH HỌA MA TRẬN KỀ =================
\begin{frame}{\faTable\ Minh họa ma trận kề}
\begin{tcolorbox}
\begin{itemize}
\item Đồ thị gồm 4 đỉnh: 1, 2, 3, 4
\item Các cạnh: (1,2), (1,3), (2,4)
\item Biểu diễn bằng ma trận kề
\item $A[i][j] = 1$ nếu có cạnh
\item $A[i][j] = 0$ nếu không có cạnh
\end{itemize}
\end{tcolorbox}

\vspace{0.3cm}

\begin{columns}

% Cột trái: đồ thị
\column{0.5\textwidth}
\centering
\begin{tikzpicture}[scale=1]

\node[circle,draw,fill=softblue!60] (1) at (0,0) {1};
\node[circle,draw,fill=softblue!60] (2) at (2,0) {2};
\node[circle,draw,fill=softpurple!60] (3) at (1,1.5) {3};
\node[circle,draw,fill=softpurple!60] (4) at (1,-1.5) {4};

\draw (1)--(2);
\draw (1)--(3);
\draw (2)--(4);

\end{tikzpicture}

% Cột phải: ma trận kề
\column{0.5\textwidth}
\begin{tcolorbox}[colframe=softpurple!70!black]
\textbf{Ma trận kề:}
\[
\begin{bmatrix}
0 & 1 & 1 & 0 \\
1 & 0 & 0 & 1 \\
1 & 0 & 0 & 0 \\
0 & 1 & 0 & 0
\end{bmatrix}
\]
\end{tcolorbox}

\end{columns}

\end{frame}
%================ SLIDE 4 =================
\begin{frame}{\faBrain\ Phân tích và ứng dụng DFS}
\begin{tcolorbox}
\begin{itemize}
\item DFS là thuật toán cơ bản trong lý thuyết đồ thị
\item Phù hợp với bài toán duyệt toàn bộ không gian trạng thái
\item Sử dụng trong tìm kiếm và trí tuệ nhân tạo
\item Là nền tảng cho nhiều thuật toán nâng cao
\item Ví dụ: Tarjan, Kosaraju
\item Áp dụng trong phân tích mạng
\item Duyệt cây và cấu trúc dữ liệu phân cấp
\item Kết hợp với BFS trong nhiều bài toán
\item Tối ưu bằng đánh dấu visited
\item Ý nghĩa trong Tin học và Toán học
\end{itemize}
\end{tcolorbox}
\end{frame}

%================ SLIDE: VÍ DỤ MINH HỌA =================
\begin{frame}{\faLightbulb\ Ví dụ minh họa DFS}
\begin{tcolorbox}
\begin{itemize}
\item Xét đồ thị $G = (V,E)$ với $V = \{1,2,3,4\}$
\item Các cạnh: (1,2), (1,3), (2,4)
\item Bắt đầu DFS từ đỉnh 1
\item Thứ tự duyệt: 1 → 2 → 4 → 3
\item Sử dụng mảng visited để tránh lặp
\end{itemize}
\end{tcolorbox}

\vspace{0.3cm}

\begin{center}
\begin{tikzpicture}[scale=1]

\node[circle,draw,fill=softblue!60] (1) at (0,0) {1};
\node[circle,draw,fill=softblue!60] (2) at (2,0) {2};
\node[circle,draw,fill=softpurple!60] (3) at (1,1.5) {3};
\node[circle,draw,fill=softpurple!60] (4) at (1,-1.5) {4};

\draw (1)--(2);
\draw (1)--(3);
\draw (2)--(4);

\end{tikzpicture}
\end{center}

\end{frame}

%================ SLIDE: CODE DFS =================
\begin{frame}{\faCode\ Ví dụ code DFS (C++)}
\begin{tcolorbox}
\begin{itemize}
\item DFS duyệt theo chiều sâu
\item Dùng đệ quy để cài đặt
\item Sử dụng danh sách kề
\item Có mảng visited để tránh lặp
\item Độ phức tạp $O(V + E)$
\end{itemize}
\end{tcolorbox}

\vspace{0.3cm}

\begin{lstlisting}
#include <iostream>
#include <vector>
using namespace std;

vector<int> adj[5];
bool visited[5];

void dfs(int u) {
    visited[u] = true;
    cout << u << " ";

    for(int v : adj[u]) {
        if(!visited[v]) {
            dfs(v);
        }
    }
}

int main() {
    adj[1] = {2, 3};
    adj[2] = {1, 4};
    adj[3] = {1};
    adj[4] = {2};

    dfs(1);
    return 0;
}
\end{lstlisting}

\end{frame}


%================ SLIDE: GIẢI THÍCH =================
\begin{frame}{\faLightbulb\ Giải thích DFS}
\begin{tcolorbox}[colframe=softpurple!70!black]
\begin{itemize}
\item Bắt đầu từ đỉnh 1
\item Đi sâu: 1 → 2 → 4
\item Sau đó quay lui (backtrack)
\item Tiếp tục duyệt đỉnh 3
\item visited giúp tránh lặp
\item Kết quả: 1 2 4 3
\end{itemize}
\end{tcolorbox}
\end{frame}
%================ SLIDE: ĐỘ PHỨC TẠP =================
\begin{frame}{\faChartLine\ Độ phức tạp}
\begin{tcolorbox}
\begin{itemize}
\item Danh sách kề: duyệt $O(V+E)$
\item Ma trận kề: duyệt $O(V^2)$
\item Kiểm tra cạnh:
\item List: $O(\deg(u))$
\item Matrix: $O(1)$
\end{itemize}
\end{tcolorbox}
\end{frame}

%================ SLIDE: ỨNG DỤNG =================
\begin{frame}{\faCogs\ Ứng dụng trong thuật toán}
\begin{tcolorbox}
\begin{itemize}
\item DFS, BFS → dùng danh sách kề
\item Dijkstra → list nhanh hơn
\item Floyd → dùng ma trận kề
\item Tối ưu phụ thuộc cấu trúc đồ thị
\item Ứng dụng trong AI, mạng
\end{itemize}
\end{tcolorbox}
\end{frame}

%================ SLIDE: ĐỒ THỊ TRỌNG SỐ =================
\begin{frame}{\faWeight\ Đồ thị có trọng số}
\begin{tcolorbox}
\begin{itemize}
\item Cạnh có giá trị trọng số
\item List: lưu (đỉnh, trọng số)
\item Matrix: lưu trực tiếp giá trị
\item Không có cạnh → $\infty$
\item Dùng trong bài toán đường đi
\end{itemize}
\end{tcolorbox}

\vspace{0.3cm}

\begin{center}
\begin{tikzpicture}[scale=1]

% Nodes
\node[circle,draw,fill=softblue!60] (1) at (0,0) {1};
\node[circle,draw,fill=softblue!60] (2) at (3,0) {2};
\node[circle,draw,fill=softpurple!60] (3) at (1.5,1) {3};
\node[circle,draw,fill=softpurple!60] (4) at (1.5,-1) {4};

% Edges (KHÔNG giao nhau)
\draw (1)-- node[above, fill=white] {2} (3);
\draw (3)-- node[right, fill=white] {5} (2);
\draw (2)-- node[below, fill=white] {3} (4);
\draw (4)-- node[left, fill=white] {4} (1);

\end{tikzpicture}
\end{center}

\end{frame}

%================ SLIDE: ĐỒ THỊ CÓ HƯỚNG =================
\begin{frame}{\faArrowRight\ Đồ thị có hướng}
\begin{tcolorbox}
\begin{itemize}
\item Cạnh có hướng $(u \rightarrow v)$
\item Ma trận: $A[i][j] \neq A[j][i]$
\item List: chỉ lưu cạnh đi ra
\item Áp dụng trong mạng, web
\item Ví dụ: đồ thị DAG
\end{itemize}
\end{tcolorbox}

\vspace{0.3cm}

\begin{center}
\begin{tikzpicture}[scale=1, >=stealth]

\node[circle,draw,fill=softblue!60] (1) at (0,0) {1};
\node[circle,draw,fill=softblue!60] (2) at (2,0) {2};
\node[circle,draw,fill=softpurple!60] (3) at (1,1.5) {3};
\node[circle,draw,fill=softpurple!60] (4) at (1,-1.5) {4};

\draw[->, thick] (1) -- (2);
\draw[->, thick] (2) -- (3);
\draw[->, thick] (3) -- (1);
\draw[->, thick] (1) -- (4);

\end{tikzpicture}
\end{center}

\end{frame}
%================ SLIDE: CHUYỂN ĐỔI =================
\begin{frame}{\faExchangeAlt\ Chuyển đổi biểu diễn}
\begin{tcolorbox}
\begin{itemize}
\item Matrix → List: duyệt từng hàng
\item List → Matrix: gán lại giá trị
\item Giữ nguyên cấu trúc đồ thị
\item Dùng trong tối ưu thuật toán
\item Phổ biến trong lập trình thi
\end{itemize}
\end{tcolorbox}
\end{frame}

%================ SLIDE: BỘ NHỚ =================
\begin{frame}{\faMemory\ Phân tích bộ nhớ}
\begin{tcolorbox}
\begin{itemize}
\item Matrix: $O(V^2)$
\item List: $O(V+E)$
\item Đồ thị thưa → dùng list
\item Đồ thị dày → dùng matrix
\item Tối ưu theo bài toán
\end{itemize}
\end{tcolorbox}
\end{frame}

%================ SLIDE: PHÉP TOÁN =================
\begin{frame}{\faCalculator\ Phép toán trên đồ thị}
\begin{tcolorbox}
\begin{itemize}
\item Tính bậc đỉnh
\item Matrix: tổng hàng/cột
\item List: số phần tử kề
\item Kiểm tra cạnh tồn tại
\item Duyệt toàn bộ đồ thị
\end{itemize}
\end{tcolorbox}
\end{frame}

%================ SLIDE: SO SÁNH =================
\begin{frame}{\faBalanceScale\ So sánh hai cách biểu diễn}
\begin{tcolorbox}
\begin{itemize}
\item List: tiết kiệm bộ nhớ
\item Matrix: truy cập nhanh
\item List: phù hợp DFS/BFS
\item Matrix: đơn giản hơn
\item Chọn theo bài toán
\end{itemize}
\end{tcolorbox}
\end{frame}
\end{document}
